% Laget automatisk av elevark/lag_elevark.py fra opplegg/kraftsensorer.html. Endringer her blir overskrevet.
\documentclass[11pt]{article}
\usepackage{elevark}
\begin{document}
\elevarkhode{Fysikk 1 · Bevegelse og krefter}{Kraftsensorer og Newtons lover}{https://fysikklab.no/opplegg/kraftsensorer.html}{Opplegget på nett}

En kraftsensor måler kraften som virker på kroken. Her er sju korte stasjoner der dere blir kjent med sensoren og samtidig ser Newtons tre lover i praksis. Skriv ned hva dere gjetter før hver måling.

\section*{Utstyr}

\begin{itemize}
\item To eller tre kraftsensorer med datalogger eller app (for eksempel Capstone eller SPARKvue), eller fjærvekter.
\item Lodd med kjent masse (100–500~g), stativ, hyssing og en vekt.
\item En vogn eller Smart Cart, gjerne med akselerasjonsmåler, og en kloss.
\end{itemize}

\section*{Stasjon 1: Nullstill og sjekk sensoren}

Nullstill sensoren uten belastning, i samme stilling som den skal brukes. Heng på lodd på 100, 200, 300, 400 og 500~g, og les av kraften. Skriv inn i tabellen.
\begin{center}\begin{tabular}{l*{5}{p{1.6cm}}}\toprule
$m$ (g) & 100 & 200 & 300 & 400 & 500 \\ \midrule
$F$ (N) & & & & & \\[4pt] \bottomrule
\end{tabular}\end{center}
 Stigningstallet til $F$ mot $m$ skal bli $g$.

\section*{Stasjon 2: Fortegn}

Dra i kroken og dytt på den. Hvilket fortegn får kraften i hvert tilfelle? Noen sensorer og programmer lar dere snu fortegnet. Avtal hva som er positivt før dere måler videre.

\section*{Stasjon 3: To sensorer som drar i hverandre}

Hekt to sensorer sammen krok mot krok. Den ene eleven holder den ene sensoren i ro, og den andre drar og rykker. Vis begge målingene i samme graf. Gjenta mens begge drar, og mens begge flytter seg.

\section*{Stasjon 4: To sensorer i serie}

Heng den ene sensoren i stativet, den andre sensoren i kroken på den første, og et lodd nederst. Nullstill begge mens de henger uten lodd. Hva viser hver av dem? Gjenta uten å nullstille den øverste med den nederste hengende på.

\section*{Stasjon 5: To sensorer i parallell}

Heng et lodd i to sensorer som holdes side om side. Hva viser hver sensor, og hva blir summen? Dra så sensorene fra hverandre, slik at snorene danner en vinkel. Hva skjer?

\section*{Stasjon 6: Jevn fart}

Dra en kloss eller vogn langs bordet med sensoren, først fra ro og så med jevn fart. Hva viser sensoren når farten er konstant? Sammenlign med kraften som trengs for å få den i gang.

\section*{Stasjon 7: Kraft og akselerasjon}

Fest sensoren i en vogn og dra med ulike, jevne krefter. Mål akselerasjonen med vognas akselerasjonsmåler, en bevegelsessensor eller videoanalyse. Skriv inn $F$ og $a$ i tabellen. Stigningstallet til $F$ mot $a$ skal bli massen til vogna.

\section*{Oppgaver}

\begin{oppgaveliste}
\oppgave \fase{test}{Mål}Stasjon 1: Hvor nær 9,81~N/kg kommer dere? Hvorfor er det viktig å nullstille sensoren i den stillingen den skal brukes?
\svar
\oppgave \fase{gjett}{Gjett}Stasjon 3: Viser sensoren til den som drar mer enn sensoren til den som holder igjen? Hva viser målingene? Hvilken av Newtons lover er dette?
\svar
\oppgave \fase{forklar}{Forklar}Stasjon 4: Hvorfor viser den øverste sensoren mer enn den nederste når den ikke er nullstilt med den andre sensoren hengende på? Hvor mye mer?
\svar
\oppgave \fase{forklar}{Forklar}Stasjon 5: Hva skjer med kraften i hver sensor når vinkelen mellom snorene blir større? Hvorfor?
\svar
\oppgave \fase{forklar}{Forklar}Stasjon 6: En elev sier at det trengs en kraft for å holde noe i bevegelse. Hva viser forsøket, og hvordan forklarer Newtons 1. lov det?
\svar
\oppgave \fase{test}{Mål}Stasjon 7: Hva blir stigningstallet til $F$ mot $a$? Sammenlign med massen dere målte med vekta.
\svar
\end{oppgaveliste}

\end{document}
